\chapter{Documenta adecuadamente el proyecto}
\label{cap:documentacion}

Este capítulo reúne la documentación que respalda el desarrollo: los registros de decisión de
diseño, los cambios implementados en esta etapa con las desviaciones declaradas respecto del
plan original, y la evidencia de los repositorios git con su convención de mensajes.

\section{Registros de decisiones de diseño}
\label{sec:registros}

El registro tiene treinta y dos decisiones: D-01 a D-21, fechadas al 2 de octubre de 2026 en
el documento de tesis, y D-22 a D-32, agregadas en esta etapa al releer el código (anexo
\ref{anexo:decisiones}). Los nueve registros siguientes agrupan las significativas, las
costosas de revertir \parencite[cap.~2, §2.4]{swebok2025}, con la fundamentación que pide
SWEBOK \parencite[cap.~3, §4.6]{swebok2025}; la razón y la evidencia de cada decisión están en
el cuadro \ref{tab:decisiones}.

\begingroup
\setbeforeparaskip{-1ex}
\paragraph{RD-1. Motores (D-11, D-12, D-25).} Para auditar el estimador y aislar la mejora,
se separan port y NLopt. Se descartan paquetes R y estimadores bayesianos
\parencite{clinton2004statistical,martin2002dynamic,shin2025l1ideal}; Fortran sigue como oráculo. Vigente.

\paragraph{RD-2. Organización (D-19, D-22--D-24, D-31).} Trazar arreglos exige estado
concentrado y búsquedas libres; GRID es configurable. Se descartan objetos por fase y un
único linaje, sin afirmar ocultamiento completo. Vigente.

\paragraph{RD-3. Configuración (D-01, D-02).} Una SVD sustituida invalidó la reproducción.
El perfil wmay/fiel fija DGESDD y aritmética estricta: se descartan sustitución automática
y \emph{fast-math}; otra configuración requiere otro sello. Vigente.

\paragraph{RD-4. Hilos (D-03, D-28).} El acumulador compartido dispersaba corridas:
suma ordenada y comparación a un hilo, sin reducciones flotantes OpenMP. Las pruebas
manuales no cierran P-01; I1 exige corrección y prueba 1/2/4/12. Vigente, corrección pendiente.

\paragraph{RD-5. Ciclos (D-07, D-32).} Sin convergencia declarada, se comparan ciclos
iguales; se descarta una parada común. El moderno controla el descenso mediante su propio
mecanismo. Vigente.

\paragraph{RD-6. Portal (D-08, D-10, D-26--D-30).} Las banderas divergían entre llamadas:
un contrato R impone aceptación, excluidos, sello y tiempo servido, sin migrar todos los
guiones. Sin sello se regenera; los hilos faltan en el resumen fiel (P-02). Vigente.

\paragraph{RD-7. Verificación (D-14, D-15).} Un error apenas movió D1 y sí la
verosimilitud \parencite{nieves2026jcc}: se exige componente, estado, mapa y falsación,
descartando solo salidas finales o umbrales de corridas que pasan. Vigente; umbral pendiente.

\paragraph{RD-8. Chile (D-16, D-17).} Periodización y D2 afectan los resultados
\parencite{nieves2026jcc}: ventana igualada e inferencia en D1, frente a ventana expansiva
y ambas dimensiones. D-16 por ratificar; D-17 vigente.

\paragraph{RD-9. Rendimiento (D-20, D-21).} La desventaja frente a wmay
\parencite{nieves2026jcc} motiva medir contra el original 2004 y contra sí mismo. Se exige
salida fiel idéntica; se descartan tolerancias y tiempos publicados como base. E3 propuesto;
P-01 primero.
\endgroup

\section{Cambios implementados y desviaciones declaradas}
\label{sec:cambios}

Los cambios de esta etapa son de diseño y de documentación; no se implementó código nuevo en
los motores, que existen desde el trabajo del artículo con el banco de verificación y el
paquete de reproducción. El cuadro \ref{tab:cambios-etapa} los registra con fecha y razón, y
el cuadro \ref{tab:desviaciones} declara las desviaciones respecto del plan original.

\footnotesize
\begin{longtable}{L{1.4cm}L{7.0cm}L{4.2cm}}
\caption{Cambios implementados en la Etapa 2.}\label{tab:cambios-etapa}\\
\toprule
Fecha & Cambio & Por qué \\
\midrule
\endfirsthead
\multicolumn{3}{l}{\footnotesize\emph{Cuadro \ref{tab:cambios-etapa}, continuación}}\\
\toprule
Fecha & Cambio & Por qué \\
\midrule
\endhead
\bottomrule
\endfoot
10 sep. & registro de tareas retirado del documento de tesis & vive en la planilla \\
24 sep. & capítulo 4, Diseño, y Anexo A, Registro de diseño & entregable de diseño \\
2 oct. & capítulo 2 completado; capítulo 4 revisado (atributos de calidad, mejora de
rendimiento, D-20, D-21); cifras corregidas & indicación del 28 sept.; auditoría de cifras \\
4 oct. & revisión adversarial independiente: 31 hallazgos, 28 corregidos y 2 en parte; RNF09 &
control de calidad \\
5 a 8 oct. & juego de auditoría de nueve vistas; figuras F1 a F4, F6, F7, A1 y A2 y cuadro T5, corregidos tras una
revisión independiente; D-22 a D-32; cuadros de atributos, niveles, construcción, errores y
perfil; este informe & indicaciones del 5 oct.; ningún pie daba la razón de una decisión \\
\end{longtable}

\begin{longtable}{L{2.6cm}L{5.3cm}L{4.7cm}}
\caption{Desviaciones respecto del plan original y su justificación.}\label{tab:desviaciones}\\
\toprule
Desviación & Qué cambió & Por qué \\
\midrule
\endfirsthead
\multicolumn{3}{l}{\footnotesize\emph{Cuadro \ref{tab:desviaciones}, continuación}}\\
\toprule
Desviación & Qué cambió & Por qué \\
\midrule
\endhead
\bottomrule
\endfoot
Entregable & informe propio por criterio & acuerdo del 5 de octubre \\
Planificación & E3 y H8 adelantadas a la Etapa 2; H1 a H9 agregadas & todo el diseño se mide
al cierre de la Etapa 2 \\
Implementación & A4 y A5 no se iniciaron & el esfuerzo fue al diseño; se reprograman \\
Requisitos & RNF09 después de la Etapa 1 & el diseño se apoya en ese atributo \\
Plazo de OE1 & mejora completa medida a inicios de diciembre & alcanzable en el curso \\
Extensión de la tesis & el cuerpo pasó de 40 a 53 páginas y bajó a 50 & se reduce hacia la
Etapa 4 \\
Commits & Etapa 1 en un commit; desde la Etapa 2, uno por pasada & el criterio lee el
historial \\
Utilidad & sin el factor $1/2$ del artículo, declarado & reproduce las tres implementaciones;
la errata queda para sus autores \\
\end{longtable}
\normalsize

\section{Repositorios git y evidencia de commits}
\label{sec:git}

El cuadro \ref{tab:repos} excluye cambios sin commit. RNF08 exige
\texttt{área(ID): verbo y qué}, razón en el cuerpo y un commit por pasada del documento.
Portal y backlog se verificaron; sus estados difieren del registro de artefactos. Diseño y reorganización siguen locales
(anexo \ref{anexo:procedencia}). Publicación y planilla quedan pendientes.

\begin{table}[htbp]
\centering
\caption{Repositorios del proyecto y su actividad, medida el 8 de octubre de 2026.}
\label{tab:repos}
\footnotesize
\begin{tabular}{L{3.4cm}L{1.75cm}L{2.1cm}L{5.0cm}}
\toprule
Repositorio & Visibilidad & Commits desde 4 sept. & Observación y ejemplo de mensaje \\
\midrule
\href{https://github.com/cowerino/dwnominate-chile}{Motores y reproducción} &
público & 0 públicos; 6 locales & rama pública \archivo{00ba24c0}; reorganización local en \archivo{ad945a91} \\
\addlinespace
\href{https://github.com/quevotan/quevotan-db/commits/fc25b1c9302cfd7b35a265357c203f21848983b6}{Investigación y backlog} &
privado, del equipo & 37 remotos, sin fusiones & \archivo{0ab30d3}: corrige formatos de semillas y documenta ajustes fallidos \\
\addlinespace
Documento de tesis e informes & privado & 8 locales & \archivo{40dae6b}: corrige diseño contra fuentes; esta revisión sigue sin commit \\
\bottomrule
\end{tabular}
\end{table}
